\chapter{Tangential Lipschitz calculus for currents}\label{sec:tangential-lipschitz-calculus}
The passage from a smooth cutoff to a Lipschitz slicing function requires convergence on the rectifiable support, which can have zero ambient volume. We make this step explicit. The source is \cite{Sim18}, Chapter 6, Lemma 4.5, pp.~182--183; the countable $C^1$ decomposition of rectifiable sets and the Euclidean area formula remain background inputs.

\begin{proposition}[Scalar coarea on a rectifiable curve]\label{prop:curve-scalar-coarea}
Let $M\subset U\subset\mathbb R^N$ be a measurable countably one-rectifiable set of locally finite $\mathcal H^1$ measure, and let $f:U\to\mathbb R$ be locally Lipschitz. Define $\nabla^Mf$ by differentiating on $C^1$ curve pieces. For every nonnegative measurable $g$,
\[
 \int_M g|\nabla^Mf|\,d\mathcal H^1
 =\int_{\mathbb R}\sum_{x\in M\cap f^{-1}(t)}g(x)\,dt.
\]
The formula is understood with measurable representatives; exceptional null subsets of $M$ contribute zero for almost every $t$. In particular the critical part $\{\nabla^Mf=0\}$ contributes no points to almost every fiber. This is the counting-measure endpoint of the scalar coarea identity used in \cite{Sim18}, Chapter 6, Lemma 4.1.
\end{proposition}
\begin{proof}
Use the countable $C^1$ decomposition to partition $M$, up to an $\mathcal H^1$-null set, into disjoint measurable pieces contained in injective regular curve charts $\gamma:I\to\mathbb R^N$. Subdivide the charts so that they and their inverses on their images are Lipschitz on the relevant relatively compact intervals. The weighted area formula, Proposition \ref{prop:weighted-area}, first gives arclength integration with factor $|\gamma'|$. The ordinary chain rule for the Lipschitz function $f\circ\gamma$ gives
\[
 |(f\circ\gamma)'|=|\nabla^Mf|\,|\gamma'|
\]
almost everywhere on each piece. Apply the one-dimensional weighted area formula to $f\circ\gamma$, inserting the indicator of that piece and the weight $g\circ\gamma$. The resulting sum counts exactly its points at level $t$, since $\gamma$ is injective. Sum over the disjoint pieces by Tonelli. The omitted null set has an $\mathcal L^1$-null image under the locally Lipschitz scalar map, by the Hausdorff-measure Lipschitz bound on a countable compact cover. Thus it affects no almost-everywhere assertion. Setting $g$ equal to the critical-set indicator proves the final assertion. This endpoint does not require applying the strict inequality $k<m$ in Theorem \ref{thm:rectifiable-coarea} with $k=m=1$.
\end{proof}

\begin{lemma}[Weak tangential convergence of ambient mollifications]\label{lem:weak-tangential-mollification}
Let $M$ be countably $m$-rectifiable with locally finite $\mathcal H^m$ measure, and let $d\mu=\theta\,d\mathcal H^m\llcorner M$ be locally finite, where $\theta\geq0$ is measurable. Suppose $h$ is locally Lipschitz in an ambient open neighborhood. On a fixed compact subset, let $h_\varepsilon$ be smooth ambient mollifications of $h$. If $v$ is a $\mu$-integrable tangent vector field supported in that compact set, then
\[
 \lim_{\varepsilon\downarrow0}\int_M v\cdot\nabla^Mh_\varepsilon\,d\mu
 =\int_M v\cdot\nabla^Mh\,d\mu.
\]
This is a weak convergence statement against fixed $L^1(\mu)$ fields. It does not assert pointwise convergence of ambient derivatives on $M$. See \cite{Sim18}, p.~182, equation (1) in the proof of Lemma 4.5.
\end{lemma}
\begin{proof}
Localize to a slightly larger compact neighborhood on which $h$ has a finite Lipschitz constant $L$. A cutoff and an extension can be used before convolution; sufficiently small convolution radii give uniform convergence to $h$ and a uniform gradient bound on the compact set being tested.

First work in a relatively compact regular $C^1$ chart $F:V\subset\mathbb R^m\to M_j$, with $DF$ and the inverse metric bounded. The functions $h_\varepsilon\circ F$ converge uniformly to $h\circ F$, while their first derivatives are uniformly bounded. For a smooth compactly supported scalar test $a$, integration by parts gives
\[
 \int_V a\,\partial_i(h_\varepsilon\circ F)
 =-\int_V(\partial_i a)(h_\varepsilon\circ F)
 \longrightarrow\int_V a\,\partial_i(h\circ F).
\]
Approximation of an $L^1$ test by smooth tests, together with the uniform derivative bound, extends this convergence to every fixed $L^1$ coefficient. Expressing $v\cdot\nabla^Mh_\varepsilon$ in chart coordinates introduces only bounded coefficients from $DF$ and the inverse metric. The factor $\theta$ and the chart area density are absorbed into the $L^1$ test; the indicator of a measurable rectifiable piece is absorbed there as well. No derivative of a merely continuous chart metric is taken.

Partition $M$ into disjoint measurable pieces covered by such charts, using the countable $C^1$ decomposition. A finite number of compact chart pieces captures all but an arbitrarily small amount of $\int|v|\,d\mu$. On that remainder both tangential gradients have the same uniform bound, so the difference of the integrals is bounded by a constant times this small tail. First pass to the limit on the finite union, then let the tail tend to zero. This proves the assertion for the original field.
\end{proof}

\begin{lemma}[Smooth scalar chain rule on a rectifiable support]\label{lem:tangential-scalar-chain}
Under the preceding rectifiability hypotheses, let $f$ be locally Lipschitz and $Q\in C^1(\mathbb R)$. Then
\[
 \nabla^M(Q\circ f)=Q'(f)\nabla^Mf
 \quad\mathcal H^m\text{-almost everywhere on }M.
\]
On each compact set where $f$ and $Q'$ are bounded, the composition is Lipschitz. This is the smooth-outer-function case of the chain rule in \cite{Sim18}, Chapter 6, p.~183, equation (4).
\end{lemma}
\begin{proof}
In each $C^1$ chart the Lipschitz function $f\circ F$ is differentiable almost everywhere by Theorem \ref{thm:rademacher}. The ordinary chain rule there gives $D(Q\circ f\circ F)=Q'(f\circ F)D(f\circ F)$. Passing from coordinate differentials to the tangential gradient proves the identity on that chart. The countable decomposition removes a single $\mathcal H^m$-null set. The mean value theorem bounds the Lipschitz constant of the composition on each relevant compact neighborhood.
\end{proof}

\begin{definition}[Tangential contraction of a rectifiable current]\label{def:tangential-current-contraction}
Let $T=\llbracket M,\theta,\tau\rrbracket$ be an integer rectifiable $m$-current of locally finite mass, $m\geq1$, and let $h$ be locally Lipschitz. Set
\[
 (T\llcorner d^Mh)(\omega)
 =\int_M\langle\tau,d^Mh\wedge\omega\rangle\theta\,d\mathcal H^m.
\]
Here $d^Mh$ can be extended to the ambient cotangent space by zero on the normal space; any extension gives the same evaluation on $\tau$. Its mass measure is bounded by $|\nabla^Mh|\,\|T\|$. This convention extends the smooth contraction in Definition \ref{def:current-restriction}; compare \cite{Sim18}, Chapter 6, equation (4.3) and the proof of Lemma 4.5.
\end{definition}

\begin{proposition}[Boundary of a Lipschitz scalar restriction]\label{prop:lipschitz-rectifiable-product-rule}
Let $T=\llbracket M,\theta,\tau\rrbracket$ be locally integer rectifiable and locally normal, with $m\geq1$, and let $h$ be locally Lipschitz. Then
\[
 \partial(T\llcorner h)=(\partial T)\llcorner h-T\llcorner d^Mh.
\]
Thus $T\llcorner h$ is locally normal, although its multiplicities need not be integers, and
\[
 \|\partial(T\llcorner h)\|
 \leq |h|\,\|\partial T\|+|\nabla^Mh|\,\|T\|.
\]
These are measure inequalities on relatively compact sets. The argument expands the Lipschitz approximation step in \cite{Sim18}, Lemma 4.5, pp.~182--183.
\end{proposition}
\begin{proof}
Fix a smooth compactly supported $(m-1)$-form $\omega$. Apply Proposition \ref{prop:current-cutoff-boundary} to ambient mollifications $h_\varepsilon$. Uniform convergence on the support of the tests gives convergence of $(T\llcorner h_\varepsilon)(d\omega)$ and $((\partial T)\llcorner h_\varepsilon)(\omega)$, since both currents have locally finite mass. The remaining term is
\[
 \int_M\langle\tau,d^Mh_\varepsilon\wedge\omega\rangle\,d\|T\|.
\]
Its integrand is the pairing of $\nabla^Mh_\varepsilon$ with a fixed bounded measurable tangent field determined by $\tau$ and $\omega$. Lemma \ref{lem:weak-tangential-mollification} gives its limit. This proves the current identity. Contraction of a unit simple $m$-vector with a covector of norm $a$ has mass norm at most $a$; the integral representation therefore gives the stated bound on every Borel set by domination of vector-measure variation. The right side is locally finite. Integer multiplicity of $T$ also ensures locally finite $\mathcal H^m$ on its nonzero multiplicity support, as required in the tangential-calculus lemmas.
\end{proof}
